\documentclass[11pt,a4paper]{article}
\usepackage[margin=0.95in]{geometry}
\usepackage{amsmath,amssymb,amsthm,mathtools}
\usepackage{enumitem}
\usepackage{xcolor}
\usepackage{titlesec}
\usepackage{mdframed}
\usepackage{hyperref}

\titlespacing*{\section}{0pt}{12pt}{6pt}
\titlespacing*{\subsection}{0pt}{8pt}{4pt}
\titleformat{\section}{\large\bfseries}{\thesection.}{0.6em}{}
\titleformat{\subsection}{\normalsize\bfseries}{\thesubsection}{0.6em}{}

\newmdenv[
  linecolor=blue!40, linewidth=0.4pt, roundcorner=2pt,
  backgroundcolor=blue!4, innerleftmargin=8pt, innerrightmargin=8pt,
  innertopmargin=5pt, innerbottommargin=5pt
]{tldr}
\newmdenv[
  linecolor=gray!40, linewidth=0.4pt, roundcorner=2pt,
  backgroundcolor=gray!5, innerleftmargin=8pt, innerrightmargin=8pt,
  innertopmargin=5pt, innerbottommargin=5pt
]{nbox}

\setlength{\parskip}{4pt}
\setlength{\abovedisplayskip}{5pt}
\setlength{\belowdisplayskip}{5pt}

\title{\vspace{-2em}\textbf{Plan: Overworked Americans take-home}\\[0.2em]
\large What I am building, in plain English, for whoever wants to follow}
\author{Alessandro Caggia}
\date{May 2026}

\begin{document}
\maketitle
\vspace{-2em}

\begin{tldr}
\textbf{The whole story in one paragraph.} Q1 says: in the data, EU
hours fell vs.\ US, and EU tax wedges rose vs.\ US -- but the
cross-sectional fit between hours and taxes is too noisy
($R^2\!=\!0.16$) for taxes to be the whole story. Q2 puts the simplest
labour-supply model (Prescott log--log) on this data and asks: what
elasticity would the model need to fit the EU--US gap? The answer is
$\sim\!0.7$ -- much larger than what the micro labour-supply
literature gives us ($\sim\!0.3$, Chetty--Guren--Manoli--Weber 2011).
So either the macro elasticity really is that big (the original
Prescott bet), or there are \emph{other things} -- institutional or
cultural -- pulling EU hours down at common $\tau$. We pursue the
second route. Q4 asks: what if Europe has hours regulation that the
US lacks, and that regulation is binding? Q5 asks: what if Europe has
unions and the US essentially does not, and unions bargain hours
toward the level workers actually prefer? Q6 asks: what if leisure in
Europe is more valuable because everyone else is also on holiday?
Q7 puts all the channels in one frame and asks how much each one
contributes to the US--EU welfare gap.
\end{tldr}

% =========================================================================
\section{The setup, in plain English}
% =========================================================================

We will use the same firm--worker pair throughout Q4--Q5--Q6, with one
small twist per question. The pair is described by two simple payoff
functions:
\[
\underbrace{V_w(h)\;=\;y\cdot h\;-\;\tfrac{\alpha}{2}h^2}_{\text{worker payoff}},
\qquad
\underbrace{V_f(h)\;=\;A\cdot h\;-\;\tfrac{\beta}{2}h^2\;-\;w\cdot h}_{\text{firm payoff}}.
\]
The worker's payoff is income minus a quadratic disutility of work.
The income piece is $y\cdot h$ where $y\!\equiv\!(1-\tau)w$ is take-home
hourly wage and $h$ is hours -- so $y\cdot h$ is the worker surplus
function the question prompt asks for, with the tax wedge built in.
The firm's payoff is revenue (concave in $h$, with diminishing returns)
minus the wage bill.

Three reference values of $h$ matter throughout the rest of the problem:
\begin{itemize}
\item $h^\ast \equiv \arg\max V_w = \dfrac{(1-\tau)w}{\alpha}$ -- the
hours \emph{the worker would choose} if free to pick.
\item $h^{\rm firm} \equiv \arg\max V_f = \dfrac{A-w}{\beta}$ -- the
hours \emph{the firm would choose} given a fixed wage.
\item $h^{\rm soc} \equiv \arg\max[V_w+V_f] = \dfrac{(1-\tau)w + (A-w)}{\alpha+\beta}$
-- the hours that maximize \emph{joint surplus} (the social planner's
choice).
\end{itemize}
At the baseline calibration ($\alpha\!=\!\beta\!=\!2$, $A\!=\!2$,
$w\!=\!1$, $\tau\!=\!0.4$) these are $h^\ast\!=\!0.30$,
$h^{\rm soc}\!=\!0.40$, $h^{\rm firm}\!=\!0.50$. The whole question is
about \emph{which of these three the equilibrium settles at}, and how
to move it.

\subsection{Why we need an institutional/cultural channel at all}

Imagine US and EU face the same wage $w$ and the same tax $\tau$. Two
clean benchmarks:
\begin{itemize}
\item Under \emph{perfect Walrasian clearing}, the worker freely picks
$h^\ast(\tau)$ in both countries. There is no US--EU hours gap at common
$\tau$.
\item Under \emph{firm-power} (the firm sets $h$ unilaterally, the
worker accepts), the firm picks $h^{\rm firm}$ in both countries. Again
no US--EU gap at common $\tau$.
\end{itemize}

Neither pure benchmark explains why EU works fewer hours than the US
at the same $\tau$. So we need \emph{something institutional or
cultural} that breaks the symmetry. That ``something'' is what Q4, Q5,
Q6 explore -- one candidate per question.

\subsection{What we deliberately leave out (footnotes-worth)}

\begin{nbox}
We do not model the \emph{extensive margin} -- whether people work at
all -- because participation depends on joint-vs-individual income tax
filing, public childcare, employment-protection legislation, and the
welfare system, all of which is a separate body of theory. The prompt
asks specifically about \emph{hours worked}, so we focus on the
intensive margin (Bick--Br\"uggemann--Fuchs-Sch\"undeln 2019 anyway show
that the extensive margin and intensive margin contribute roughly
equally to the EU--US $H/N$ gap, but each has its own machinery).

We also do not model \emph{contract stickiness} -- the fact that in
many European countries, weekly hours per worker are set in a sectoral
collective contract that is not renegotiated for years even when
preferences or productivity change. This is a real friction documented
in Cahuc--Carcillo (2014, IZA WP) and the EU Labour Force Survey, but
it is observationally equivalent (in our static framework) to a
binding hours cap, and is therefore captured implicitly through the
cap mechanism of Q4.

We also do not model \emph{sectoral composition}. Larger public
sectors in Europe imply more workers on contracts that are sticky for
political-economy reasons unrelated to firm-profit maximisation;
public-sector hours decisions are governed by a different mechanism
(collective bargaining over service provision). Our results should be
read as applying to the private-sector portion of the labour market.
\end{nbox}

% =========================================================================
\section{Q1 -- the data inconsistency}
% =========================================================================

We open with three panels of data: (a) hours per worker, percent
deviation from US, 1970--2024; (b) effective labour wedge $\tau_h$,
percent deviation from US, 1970--2007 (McDaniel 2011); (c) composite
tax wedge $\tau$ vs.\ hours per worker, 2023 cross-section across 19
OECD countries.

Panels (a) and (b) tell a coherent story: hours fell against the US
after 1970, wedges rose against the US, and the timing matches roughly.
Panel (c) is where things get interesting. The OLS slope is signed
correctly ($-9.9$ hours per percentage point of wedge), but the
$R^2\!=\!0.16$ and the residual standard deviation is $\sim\!158$
hours -- about one-third of the entire DE--US hours gap. Italy and
Greece have lower wedges than France yet work \emph{more} hours per
worker; Belgium, Germany, Austria, France span $\sim\!200$ hours within
a 6\,pp wedge band.

The cross-section of taxes vs.\ hours is too noisy for taxes to be the
whole story. Either the within-country time-series response of hours
to taxes is much larger than the cross-section delivers (the
\emph{Prescott--CGM puzzle} we go after in Q2), or the cross-section
is corrupted by country-specific institutional or cultural frictions
that the tax wedge does not capture (the \emph{Q4--Q6 candidates}).

% =========================================================================
\section{Q2 -- the Prescott model and the elasticity puzzle}
% =========================================================================

The simplest model of labour supply under taxation is Prescott's
(2004): a representative household with log--log preferences,
$u(c,h)\!=\!\log c + \alpha\log(1-h)$, faces a wage $w$, a labour tax,
a consumption tax, and a balanced lump-sum rebate. The first thing the
model gives us is a clean composite tax wedge $1-\tau\!\equiv\!(1-\tau_l)/(1+\tau_c)$
that absorbs both the labour and the consumption tax (since workers
supply labour only to consume, the consumption tax is isomorphic to a
labour tax). The first-order condition then gives the worker's
preferred hours as a closed-form function of $\tau$:
\[
h^\ast(\tau)\;=\;\frac{1-\tau}{\alpha + 1-\tau},\qquad
\varepsilon^M_{h,1-\tau}\;=\;\frac{\alpha}{\alpha+1-\tau}.
\]

\subsection{The calibrated model fits, but at the wrong elasticity}

Anchor $h^{US}\!=\!0.33$ at $\tau^{US}\!=\!0.35$. Inverting the formula
gives $\alpha\!=\!1.32$. With this single parameter pinned down, the
model's prediction at $\tau^{EU}\!=\!0.56$ is $h^{EU}\!=\!0.44/1.76\!=\!0.250$
-- and the data point is $0.247$. \emph{The model essentially fits the
cross-section.}

So why is there a puzzle? Because the same calibration pins down the
labour-supply elasticity at $\sim 0.67$--$0.75$ along the supply curve,
and the micro labour-supply literature (Chetty--Guren--Manoli--Weber
2011, summarising the quasi-experimental evidence) reports a Hicksian
intensive elasticity of about $0.30$ for prime-age workers. \emph{The
calibrated $\alpha$ implies an elasticity that is roughly $2.5\times$
the micro estimate.}

The puzzle is therefore not that the model under-explains the data,
but that the calibrated parameter is inconsistent with quasi-experimental
evidence on individual labour supply. Either:
\begin{itemize}
\item the macro elasticity really \emph{is} that large, and the micro
literature is missing something (life-cycle aggregation,
Rogerson--Wallenius 2009; tax progressivity, HSV 2017; participation
elasticity, BFL 2018 -- all genuine modern channels that close part of
the gap);
\item or there are \emph{other channels} that produce the cross-country
hours gap at low individual elasticity, so the macro $\alpha$ is
absorbing things that should be modelled separately as institutions or
culture.
\end{itemize}

\subsection{The decomposition that motivates Q4--Q7}

Pin $\alpha$ at the micro-consistent value
$\alpha^{\rm micro}\!=\!0.279$ (the value that delivers
$\varepsilon^M\!=\!0.30$ at the US point). Now the Prescott model
\emph{alone} predicts US hours of $0.700$ and EU hours of $0.612$ -- both
much higher than observed, in absolute terms, but the cross-country log
gap is only $-13.4\%$. The data shows a $-29.0\%$ log gap. So at the
micro elasticity, taxes alone explain about \emph{46\%} of the
cross-country hours gap; the remaining \emph{54\%} -- a 15.6
percentage-point log residual -- must come from somewhere outside the
Prescott mechanism.

This is the load-bearing observation. The institutional and cultural
channels of Q4--Q6 are not just nice-to-have additions to a story that
Prescott already tells; they are quantitatively necessary at any
plausible labour-supply elasticity. \emph{This is what motivates Q4,
Q5, Q6 as serious candidates rather than rhetorical ones.}

% =========================================================================
\section{Q4 -- hours regulation as a quantity instrument}
% =========================================================================

The first candidate is straightforward: \emph{governments in Europe
constrain hours, while the US essentially does not.} The 35-hour week
in France, the EU Working Time Directive, sectoral hour caps in
Germany and Italy -- all are quantity instruments that, if binding,
keep equilibrium hours below what they would otherwise be.

\subsection{The model: a cap on top of the firm-power baseline}

We use the firm-power baseline as our reference point: in the absence
of regulation, the firm picks $h^{\rm firm}\!=\!(A-w)/\beta\!=\!0.5$
and the worker accepts (think of this as the result of a long-term
contract that is sticky in the EU; in the US the same outcome arises
from a thinner labour market). The government imposes a cap $\bar h$.
The constrained equilibrium is
\[
h^{eq}(\bar h)\;=\;\min\bigl\{\bar h,\;h^{\rm firm}\bigr\}.
\]
If the cap is above $h^{\rm firm}$ it is non-binding; if below,
equilibrium hours equal the cap.

\subsection{The intuition: it's an externality problem}

The welfare condition for the cap is exactly the Pigouvian-tax
analogue from the externalities chapter. Joint surplus
$W(h)\!=\!V_w+V_f$ is concave with peak at $h^{\rm soc}\!=\!0.40$. So
moving from $h^{\rm firm}\!=\!0.50$ down toward $h^{\rm soc}\!=\!0.40$
\emph{increases} joint welfare -- the firm is producing too many hours
relative to the joint optimum, and pulling down a little corrects the
distortion. But moving \emph{past} $h^{\rm soc}$ -- e.g., to a cap at
$h\!=\!0.25$ -- now overshoots in the other direction: too few hours,
joint surplus falls below where it was at the firm-power baseline.

The clean result: the cap is welfare-improving \emph{iff}
$\bar h\!\in\!(2h^{\rm soc}-h^{\rm firm},\,h^{\rm firm})$, which at our
$\alpha=\beta$ baseline simplifies to
\[
\boxed{\;\bar h\;\in\;(h^\ast,\;h^{\rm firm})\;\text{ for the cap to be
welfare-improving}\;}
\]
Equivalently: any cap that pulls hours down toward $h^{\rm soc}$ helps;
any cap that pushes hours below the worker's individual optimum $h^\ast$
hurts (the worker is now being forced to work \emph{less} than they
themselves would prefer at this wage).

\subsection{What this says about US vs.\ EU}

In our framework, the question is whether the EU cap binds and the US
cap does not. The Aubry 35-hour reform in France maps to roughly
$\bar h\!\approx\!0.875\cdot h^{\rm firm}\!=\!0.4375$ -- comfortably
inside the welfare-improving band. The empirical evidence
(Estev\~ao--S\'a 2008; Chemin--Wasmer 2009) confirms a $\sim\!4\%$
hours decline in covered French sectors with near-zero employment
effect: a cap that bound, that worked, and that did not overshoot.

The US has the FLSA 40-hour overtime threshold but it rarely binds at
the median margin -- so for our purposes the US is at $h^{\rm firm}$
and the EU is at $\bar h$. This is one source of the EU--US hours gap
at common $\tau$.

\subsection{Why the cap can be the wrong tool}

There is a real cost to the cap that the worked example glosses over:
\emph{workers are heterogeneous in $\alpha$}. A single national cap
helps high-$\alpha$ (leisure-loving) workers, whose individual $h^\ast$
is far below $h^{\rm firm}$ anyway, and hurts low-$\alpha$
(income-loving) workers, whose individual $h^\ast$ may be above
$\bar h$. The cap therefore redistributes from low-$\alpha$ workers to
high-$\alpha$ workers, which may or may not be what a planner wants. We
flag this in the body but do not model it explicitly.

% =========================================================================
\section{Q5 -- a single union (the heart of the modelling story)}
% =========================================================================

The second candidate is more elegant. Caps are blunt: a government
bureaucrat picks a single number and applies it to everyone, with no
guarantee that it sits in the welfare-improving band. \emph{What if
the constraint on hours could emerge endogenously from the interaction
between workers and firms?} That is what unions do.

\subsection{The setup: one union, all workers inside, dues share $s$}

We model the simplest possible union: a single industry-wide union, all
workers are inside, each pays a fixed share $s\!\in\!(0,1)$ of net
labour income to the union. The union's payoff is $s\cdot V_w(h)$, so
\[
\arg\max_h\;\bigl[s\cdot V_w(h)\bigr]\;=\;\arg\max_h V_w(h)\;=\;h^\ast.
\]
The dues share scales the union's objective but does not move what the
union wants. \emph{This is the central modelling point}: the union and
the workers are perfectly aligned, no agency problem. (Compare a union
that maximises dues revenue $s\cdot wh$; that union would prefer
$h^{\rm firm}$ -- dues revenue is increasing in hours -- and pull in
the wrong direction.)

\subsection{Why this is conceptually better than a cap}

The cap of Q4 is set by a government planner who must guess
$h^{\rm soc}$ correctly. If the planner overshoots, welfare falls. The
union of Q5 \emph{endogenously} pulls hours toward $h^\ast$, which is
inside the welfare-improving band by construction (since $h^\ast$ is
\emph{at} the worker-side boundary of that band). This is the same
intuition that makes \emph{Coase} better than \emph{Pigou} when the
parties can negotiate: let the bargaining endogenously select the
right level, rather than have the government impose a quantity ex ante.

The union therefore has the potential to fix the firm-power distortion
without requiring the planner to know $h^{\rm soc}$.

\subsection{The variable that does the work: bargaining strength $\eta$}

The actual hours outcome depends on \emph{how strongly} the union can
bargain. Let $\eta\!\in\![0,1]$ be the worker's bargaining weight in
the firm--union negotiation. At $\eta\!=\!0$ (no union influence) the
firm picks $h^{\rm firm}$. At $\eta\!=\!1$ (full union dominance) the
union pulls hours all the way to $h^\ast$. For intermediate $\eta$, the
bargained outcome interpolates:
\[
h(\eta)\;=\;\frac{(1-\eta)(A-w) + \eta(1-\tau)w}{(1-\eta)\beta + \eta\alpha}.
\]
At our symmetric baseline ($\alpha\!=\!\beta\!=\!2$) this is linear in
$\eta$.

\subsection{The optimality condition for the union to be welfare-improving}

Joint surplus $W(\eta)$ is concave in $\eta$ with peak at the
$\eta^{\rm soc}$ where $h(\eta^{\rm soc})\!=\!h^{\rm soc}$. At our
baseline this is $\eta^{\rm soc}\!=\!1/2$ -- so the joint optimum is
hit when the union has \emph{half} the bargaining weight. The
welfare-improving region is
\[
\boxed{\;\eta\;\in\;(0,\;2\eta^{\rm soc})\;\text{ for the union bargain
to be welfare-improving}\;}
\]
which at baseline simplifies to $\eta\!\in\!(0,\,1)$. \emph{A
moderately strong union helps; a maximally strong union pulls hours
all the way to $h^\ast$, which is below $h^{\rm soc}$, and total
surplus falls back to the no-union level.}

This is the core finding for Q5: \textbf{the welfare effect of a union
is non-monotone in its strength}. A toothless union does nothing; an
all-powerful union over-corrects; a moderate union sits at the joint
optimum. The story closely mirrors Q4: over-shoot and under-shoot in
either direction destroy welfare.

\subsection{What this says about US vs.\ EU}

In the US, bargaining-coverage rates are around $10\%$. In Italy they
are around $80\%$, in France $\sim\!95\%$, Sweden $\sim\!90\%$. The
order-of-magnitude difference in coverage maps to a several-fold
difference in $\eta$. So the model predicts:
\begin{itemize}
\item \textbf{US:} low $\eta$, equilibrium hours close to $h^{\rm firm}$
(no union channel doing work);
\item \textbf{EU:} moderate $\eta$, equilibrium hours pulled toward
$h^\ast$, ideally near $h^{\rm soc}$ where the welfare gain is maximised.
\end{itemize}
The US--EU hours gap at common $\tau$ is then a difference in $\eta$,
not in $\alpha$ or $\tau$.

\subsection{The honest caveat: what happens under canonical right-to-manage}

The above is the cleanest pedagogical version. The canonical labour
economics literature (Nickell 1982; Layard--Nickell--Jackman 1991) uses
\emph{right-to-manage} (RTM): the union and firm bargain over the
\emph{wage} $w$, and the firm then picks $h$ on its labour-demand
curve. Under RTM, the union still pulls outcomes toward the worker
side, but it does so by raising $w$, which has a side effect: the tax
wedge $\tau wh$ rises, and that wedge is dead-weight (lost from
$V_w+V_f$). Once you account for this, the canonical result is that
$W$ falls monotonically in $\eta$ -- there is no welfare-improving
region.

The two readings (sticky-$w$ Coasian intuition vs.\ flexible-$w$
canonical RTM) give different welfare conclusions because they make
different assumptions about what is bargained. The sticky-$w$ reading
is the cleanest pedagogical version and probably the closest to the
``Coase beats Pigou'' intuition the prompt is fishing for. The
canonical RTM reading is what the cross-country macro-labour
literature actually uses. We present both in the body of the take-home
and let the reader decide.

% =========================================================================
\section{Q6 -- culture as a leisure bonus (enjoy life more = higher welfare)}
% =========================================================================

The third candidate is the most direct: \emph{Europeans enjoy leisure
more than Americans do.} This shows up empirically -- subjective
wellbeing surveys (Stevenson--Wolfers 2008, BPEA Spring), time-use
studies, and revealed preferences in cross-country labour-supply data
all point to a higher marginal value of leisure on the European side.

\subsection{Two ways to write ``culture means people enjoy life more''}

There is a real modelling choice here that matters for welfare. Two
options:

\textbf{(i) Higher disutility of work.} Increase $\alpha$ for the EU
worker, so working hurts them more. This delivers $h^\ast$ down -- but
also delivers $V_w$ \emph{down} at every $h$. This is the ``Europeans
suffer more from work'' reading. \emph{Welfare falls.}

\textbf{(ii) Higher utility of leisure.} Add an additive bonus
$\theta(1-h)$ to the worker's payoff. Each unit of leisure gives extra
utility on top of the avoided disutility of work. This delivers
$h^\ast$ down AND $V_w$ \emph{up}. This is the ``Europeans enjoy
leisure more'' reading. \emph{Welfare rises.}

The user's intuition is (ii), and it is also the right one
structurally. ``Enjoying life more'' should be welfare-improving by
construction. We adopt formulation (ii).

This is the standard preference-for-leisure parameter in the
cross-country labour-supply literature: Bowles--Park (2005, EJ 115)
on emulation and work hours; Olovsson (2009, IER 50(1)) on
Europe vs.\ US; Boppart--Krusell (2020, JPE 128(1)) for the
balanced-growth-compatible preference shift framework; Bell--Freeman
(2001, JOLE 19(1)) on relative-consumption preferences.

\subsection{The model}

Worker's payoff with cultural leisure-premium $\theta^c\!\geq\!0$:
\[
V_w(h;\theta^c)\;=\;(1-\tau)\,w\,h\;-\;\tfrac{\alpha}{2}\,h^2\;+\;\theta^c\,(1-h),
\qquad
\theta^{EU}\;>\;\theta^{US}\!=\!0.
\]
The first two terms are the standard quasi-linear payoff. The third
term is the cultural leisure bonus: each unit of leisure $1-h$
delivers $\theta^c$ of extra utility. Same $\alpha$ in both countries;
firms unchanged.

The worker's first-best is then:
\[
\boxed{\;h^\ast(\tau,\theta^c)\;=\;\frac{(1-\tau)\,w-\theta^c}{\alpha},
\qquad \frac{\partial h^\ast}{\partial \theta^c}\;<\;0.\;}
\]
The worker's payoff at the first-best is
$V_w^\ast\!=\![(1-\tau)w-\theta^c]^2/(2\alpha) + \theta^c$, and
$dV_w^\ast/d\theta^c\!=\!1-h^\ast\!>\!0$: \emph{higher cultural
premium $\Rightarrow$ higher worker welfare at the first-best}, as
intuition demands.

\subsection{Calibration}

We pick $\theta^{US}\!=\!0$ and $\theta^{EU}\!=\!0.156$. The latter is
calibrated so that
\[
h^\ast_{EU}\;=\;\frac{0.6-0.156}{2}\;=\;0.222,
\]
matching the empirical EU level used in \S 2's micro-elasticity
decomposition. The $0.156$ figure is roughly $26\%$ of the net hourly
wage $(1-\tau)w$ -- each EU hour of leisure carries an extra utility
worth about a quarter of a market hour, a magnitude consistent with
Stevenson--Wolfers (2008).

\subsection{Welfare implication: EU is welfare-DOMINANT}

Plug in the numbers. At each country's first-best:
\[
\text{US:}\;V_w^\ast\!=\!0.090,\;V_f^\ast\!=\!0.210,\;W^{US}\!=\!0.300.
\]
\[
\text{EU:}\;V_w^\ast\!=\!0.205,\;V_f^\ast\!=\!0.173,\;W^{EU}\!=\!0.378.
\]
The EU worker is strictly better off ($V_w$ rises from $0.090$ to
$0.205$, even though $h$ falls from $0.300$ to $0.222$): the leisure
premium adds direct utility. The EU firm is worse off (lower hours
worked). Joint surplus rises ($\Delta W\!=\!+0.078$) because the
worker's gain dominates the firm's loss. \emph{The cultural premium
is a free utility source from the planner's perspective.}

\subsection{Why $\theta^{EU}>\theta^{US}$? Three structural candidates}

\textbf{(a) AGS (2005) historical channel.} Post-WWII union-driven
shorter-week reforms in Europe shifted social norms toward a
``leisure-is-valuable'' equilibrium, and over time those norms
intensified leisure preferences (Bisin--Verdier 2001 give the
cultural-transmission micro-foundation). Under this reading,
\emph{the union channel of \S 5 is the upstream cause of the cultural
channel of \S 6}: Q5 unions are the proximate coordination device,
Q6 leisure preferences are the long-run residue.

\textbf{(b) Bowles (1998) endogenous preferences.} Markets shape
preferences over generations; long histories of different
labour-market institutions produce different aggregate preferences for
leisure.

\textbf{(c) Stevenson--Wolfers (2008) subjective-wellbeing residual.}
Europeans report higher subjective wellbeing per hour of leisure than
Americans -- the direct empirical anchor for $\theta^{EU}>\theta^{US}$.

\subsection{Discriminating evidence}

The cleanest test of preference-shift vs.\ coordination/spillover
explanations is whether immigrant workers from low-$h$ countries
import home-country preferences (preference shift) or assimilate to
host-country institutions (coordination). Empirical surveys
(Fern\'andez--Fogli 2009; Algan--Cahuc 2010) find evidence for both,
with second-generation immigrants partly retaining home-country
preferences but converging over generations. The leisure-bonus
specification is consistent with this evidence.

% =========================================================================
\section{Q7 -- the welfare comparison}
% =========================================================================

Now we put everything together. Take the EU equilibrium -- AGS
coordination on $h_L$ plus union bargaining at some $\eta^{EU}\!>\!0$
-- and ask: how much of the EU--US welfare gap comes from each of the
two channels (culture, union)?

\subsection{The four cases}

At common $\tau$ and common preferences, the EU can sit below the US
for one of four reasons, each with a distinct welfare answer:

\begin{itemize}
\item[\textbf{(I)}] US at $h^{\rm firm}$, EU at $h^{\rm soc}$ via cap or
moderate union: \emph{EU dominates}, by exactly the Q4--Q5 mechanism.
\item[\textbf{(II)}] US at $h^{\rm firm}$, EU below $h^\ast$ via overshoot
cap or excessive union: \emph{EU welfare-decreasing}.
\item[\textbf{(III)}] Same institutions, $\theta^{EU}\!>\!\theta^{US}$:
each country at its own first-best, but the EU is welfare-dominant
because the cultural premium adds free utility from leisure.
\item[\textbf{(IV)}] \textbf{Compounded.} Both the institutional bargain
and the cultural parameter differ -- the empirically realistic case.
\end{itemize}

\subsection{The counterfactual decomposition}

We run case (IV) explicitly. Starting from the US baseline
($\theta^{US}\!=\!0$, $\eta^{US}\!=\!0$ giving $w\!=\!1$, $h\!=\!0.5$,
$W^{US}\!=\!0.300$) and shutting down each channel in sequence, we get:

\begin{itemize}
\item \emph{Apply the EU cultural shift alone}
(raise $\theta$ to $\theta^{EU}\!=\!0.156$, keep $\eta\!=\!0$): the
firm-power IR-binding wage falls to $w\!=\!0.879$ (because the worker
accepts lower pay when leisure has direct utility), hours rise to
$h\!=\!0.561$, but the worker's $V_w$ stays at the IR baseline; firm
captures more profit. Total $W$ rises to $0.364$, $\Delta W\!=\!+0.064$.
\item \emph{Add the EU union on top} (raise $\eta$ to $0.5$ at
$\theta^{EU}$): the union pulls wages back up to $w\!=\!1.172$, hours
fall to $h\!=\!0.414$. Worker captures more surplus ($V_w$ rises),
firm less. Total $W$ rises further to $0.382$, $\Delta W\!=\!+0.018$.
\end{itemize}
Net EU--US gap: $\Delta W\!=\!+0.082$ -- \emph{EU is welfare-dominant.}

The reverse ordering (union first, then culture):
\begin{itemize}
\item Apply union alone (raise $\eta$ to $0.5$ at $\theta^{US}\!=\!0$):
$W$ falls to $0.283$, $\Delta W\!=\!-0.017$ (canonical RTM result --
union extracts wage rents but creates tax-wedge dead-weight).
\item Add culture (raise $\theta$ to $\theta^{EU}$ at $\eta\!=\!0.5$):
$W$ rises sharply to $0.382$, $\Delta W\!=\!+0.100$.
\end{itemize}
Averaging the two orderings (Shapley decomposition) gives the
path-independent contributions:
\[
\Delta W^{\rm Shapley}_{\rm culture}\;=\;+0.082,
\qquad
\Delta W^{\rm Shapley}_{\rm union}\;\approx\;0.
\]
\emph{The cultural channel does essentially all of the work.} The
union channel is welfare-improving in path A ($+0.018$, because under
EU culture the firm-power baseline over-extracts hours and the union
pulls them back) and welfare-reducing in path B ($-0.017$, the
canonical RTM result). On average across paths, it nets to roughly
zero.

\subsection{The reading}

This is the substantive answer. The cultural channel does the heavy
lifting: when Europeans value leisure more, they enjoy life more, and
joint welfare rises. The union channel is small and ambiguous in
sign on average. The dominant force in the EU--US labour-market
welfare comparison is the cultural premium on leisure, not the
bargaining institution.

This contradicts a naive reading of Q5 in isolation (where the canonical
RTM model gives unions a welfare-reducing role). Why? Because under the
EU cultural premium, firm-power equilibrium over-extracts hours: the
worker would prefer fewer hours at any given wage (higher
$\theta\!\Rightarrow\!$ lower $h^\ast$), so when the firm sets
$h\!=\!h^d(w)$ at the IR-binding wage, $h$ is too high relative to the
worker's preference. The union pulls $h$ back down toward $h^\ast$,
which improves welfare. The Q5-only ``unions destroy welfare via
tax-wedge dead-weight'' result still applies in path B (no cultural
shift), but path A activates a counter-mechanism.

\subsection{The honest welfare-units caveat}

Comparing $V_w + V_f$ across countries with different $\theta^c$ uses
the raw simulation units. The leisure-bonus specification holds
$\alpha$ identical and only adds an additive cultural premium, so the
$V_w$ comparison is on a common cardinalisation (same monetary
numeraire). Mapping to a money-metric (compensating variation in
consumption-equivalent units, the Jones--Klenow 2016 standard) leaves
the sign and order of magnitude unchanged.

\subsection{The bigger picture}

Jones--Klenow (2016) report a US--EU welfare gap of $-10\%$ to $-25\%$
in consumption-equivalent units across all channels (consumption,
leisure, life expectancy, inequality). Of these, mortality and aggregate
consumption levels carry the bulk -- Becker--Philipson--Soares (2005)
and Hall--Jones (2007) calibrate a value-of-life such that the
US--EU life-expectancy gap of $\sim 4$ years contributes a welfare
\emph{gain} for Europe of roughly $+20\%$. The labour-market block we
analyse here is therefore sign-determinate but quantitatively
second-order to demographic, TFP, and subjective-wellbeing channels
that this analysis does not address. We are answering ``how much of
the labour-market gap can each institutional channel explain?'',
not ``is Europe better off than the US overall?''.

% =========================================================================
\section{Where things stand}
% =========================================================================

The take-home is now four blocks (Q1--Q2--Q4--Q5--Q6--Q7) that share a
common firm-worker framework and answer one question each:
\begin{itemize}
\item Q1 says the data is suggestive but noisy.
\item Q2 says the simplest model fits at an implausibly large
elasticity, so we need other channels to explain at least half the
cross-country gap.
\item Q4 explores one such channel (hours regulation) and derives the
welfare-improving region for a cap.
\item Q5 explores another (a single industry-wide union) and derives
the welfare-improving region for the bargaining strength $\eta$.
\item Q6 explores a third (a leisure bonus that adds direct utility
from leisure, $\theta^{EU}\!>\!\theta^{US}\!=\!0$, after Bowles--Park
2005 / Olovsson 2009 / Boppart--Krusell 2020 / Stevenson--Wolfers
2008) and shows that ``Europeans enjoy life more'' is genuinely
welfare-improving by construction.
\item Q7 combines them and reports the welfare gap of EU vs.\ US under
the canonical model: \emph{$\Delta W\!=\!+0.082$ (EU welfare-dominant)},
with the cultural channel doing essentially all of the work and the
union channel netting to zero in the Shapley average.
\end{itemize}

\bigskip
\noindent\emph{Source files:} \texttt{problem1.tex} (the submission),
\texttt{research/build\_figure\_panel.py} (Q1 figure),
\texttt{research/sim\_models\_v2.py} (Q2 decomposition, Q5 RTM, Q6
leisure-bonus culture, Q7 layered counterfactual).

\end{document}
